%!TeX root=../sensetop.tex
\chapter[Chapter \thechapter]{}
\lettrine[lraise=0.3]{E}{dward} remained a week at the cottage; he was earnestly pressed by Mrs Dashwood to stay longer; but, as if he were bent only on self-mortification, he seemed resolved to be gone when his enjoyment among his friends was at the height. His spirits, during the last two or three days, though still very unequal, were greatly improved—he grew more and more partial to the house and environs—never spoke of going away without a sigh—declared his time to be wholly disengaged—even doubted to what place he should go when he left them—but still, go he must. Never had any week passed so quickly—he could hardly believe it to be gone. He said so repeatedly; other things he said too, which marked the turn of his feelings and gave the lie to his actions. He had no pleasure at Norland; he detested being in town; but either to Norland or London, he must go. He valued their kindness beyond any thing, and his greatest happiness was in being with them. Yet, he must leave them at the end of a week, in spite of their wishes and his own, and without any restraint on his time.

Elinor placed all that was astonishing in this way of acting to his mother's account; and it was happy for her that he had a mother whose character was so imperfectly known to her, as to be the general excuse for every thing strange on the part of her son. Disappointed, however, and vexed as she was, and sometimes displeased with his uncertain behaviour to herself, she was very well disposed on the whole to regard his actions with all the candid allowances and generous qualifications, which had been rather more painfully extorted from her, for Willoughby's service, by her mother. His want of spirits, of openness, and of consistency, were most usually attributed to his want of independence, and his better knowledge of Mrs Ferrars's disposition and designs. The shortness of his visit, the steadiness of his purpose in leaving them, originated in the same fettered inclination, the same inevitable necessity of temporizing with his mother. The old well-established grievance of duty against will, parent against child, was the cause of all. She would have been glad to know when these difficulties were to cease, this opposition was to yield,—when Mrs Ferrars would be reformed, and her son be at liberty to be happy. But from such vain wishes she was forced to turn for comfort to the renewal of her confidence in Edward's affection, to the remembrance of every mark of regard in look or word which fell from him while at Barton, and above all to that flattering proof of it which he constantly wore round his finger.

<I think, Edward,> said Mrs Dashwood, as they were at breakfast the last morning, <you would be a happier man if you had any profession to engage your time and give an interest to your plans and actions. Some inconvenience to your friends, indeed, might result from it—you would not be able to give them so much of your time. But (with a smile) you would be materially benefited in one particular at least—you would know where to go when you left them.>

<I do assure you,> he replied, <that I have long thought on this point, as you think now. It has been, and is, and probably will always be a heavy misfortune to me, that I have had no necessary business to engage me, no profession to give me employment, or afford me any thing like independence. But unfortunately my own nicety, and the nicety of my friends, have made me what I am, an idle, helpless being. We never could agree in our choice of a profession. I always preferred the church, as I still do. But that was not smart enough for my family. They recommended the army. That was a great deal too smart for me. The law was allowed to be genteel enough; many young men, who had chambers in the Temple, made a very good appearance in the first circles, and drove about town in very knowing gigs. But I had no inclination for the law, even in this less abstruse study of it, which my family approved. As for the navy, it had fashion on its side, but I was too old when the subject was first started to enter it—and, at length, as there was no necessity for my having any profession at all, as I might be as dashing and expensive without a red coat on my back as with one, idleness was pronounced on the whole to be most advantageous and honourable, and a young man of eighteen is not in general so earnestly bent on being busy as to resist the solicitations of his friends to do nothing. I was therefore entered at Oxford and have been properly idle ever since.>

<The consequence of which, I suppose, will be,> said Mrs Dashwood, <since leisure has not promoted your own happiness, that your sons will be brought up to as many pursuits, employments, professions, and trades as Columella's.>

<They will be brought up,> said he, in a serious accent, <to be as unlike myself as is possible. In feeling, in action, in condition, in every thing.>

<Come, come; this is all an effusion of immediate want of spirits, Edward. You are in a melancholy humour, and fancy that any one unlike yourself must be happy. But remember that the pain of parting from friends will be felt by every body at times, whatever be their education or state. Know your own happiness. You want nothing but patience—or give it a more fascinating name, call it hope. Your mother will secure to you, in time, that independence you are so anxious for; it is her duty, and it will, it must ere long become her happiness to prevent your whole youth from being wasted in discontent. How much may not a few months do?>

<I think,> replied Edward, <that I may defy many months to produce any good to me.>

This desponding turn of mind, though it could not be communicated to Mrs Dashwood, gave additional pain to them all in the parting, which shortly took place, and left an uncomfortable impression on Elinor's feelings especially, which required some trouble and time to subdue. But as it was her determination to subdue it, and to prevent herself from appearing to suffer more than what all her family suffered on his going away, she did not adopt the method so judiciously employed by Marianne, on a similar occasion, to augment and fix her sorrow, by seeking silence, solitude and idleness. Their means were as different as their objects, and equally suited to the advancement of each.

Elinor sat down to her drawing-table as soon as he was out of the house, busily employed herself the whole day, neither sought nor avoided the mention of his name, appeared to interest herself almost as much as ever in the general concerns of the family, and if, by this conduct, she did not lessen her own grief, it was at least prevented from unnecessary increase, and her mother and sisters were spared much solicitude on her account.

Such behaviour as this, so exactly the reverse of her own, appeared no more meritorious to Marianne, than her own had seemed faulty to her. The business of self-command she settled very easily;—with strong affections it was impossible, with calm ones it could have no merit. That her sister's affections \textit{were} calm, she dared not deny, though she blushed to acknowledge it; and of the strength of her own, she gave a very striking proof, by still loving and respecting that sister, in spite of this mortifying conviction.

Without shutting herself up from her family, or leaving the house in determined solitude to avoid them, or lying awake the whole night to indulge meditation, Elinor found every day afforded her leisure enough to think of Edward, and of Edward's behaviour, in every possible variety which the different state of her spirits at different times could produce,—with tenderness, pity, approbation, censure, and doubt. There were moments in abundance, when, if not by the absence of her mother and sisters, at least by the nature of their employments, conversation was forbidden among them, and every effect of solitude was produced. Her mind was inevitably at liberty; her thoughts could not be chained elsewhere; and the past and the future, on a subject so interesting, must be before her, must force her attention, and engross her memory, her reflection, and her fancy.

From a reverie of this kind, as she sat at her drawing-table, she was roused one morning, soon after Edward's leaving them, by the arrival of company. She happened to be quite alone. The closing of the little gate, at the entrance of the green court in front of the house, drew her eyes to the window, and she saw a large party walking up to the door. Amongst them were Sir John and Lady Middleton and Mrs Jennings, but there were two others, a gentleman and lady, who were quite unknown to her. She was sitting near the window, and as soon as Sir John perceived her, he left the rest of the party to the ceremony of knocking at the door, and stepping across the turf, obliged her to open the casement to speak to him, though the space was so short between the door and the window, as to make it hardly possible to speak at one without being heard at the other.

<Well,> said he, <we have brought you some strangers. How do you like them?>

<Hush! they will hear you.>

<Never mind if they do. It is only the Palmers. Charlotte is very pretty, I can tell you. You may see her if you look this way.>

As Elinor was certain of seeing her in a couple of minutes, without taking that liberty, she begged to be excused.

<Where is Marianne? Has she run away because we are come? I see her instrument is open.>

<She is walking, I believe.>

They were now joined by Mrs Jennings, who had not patience enough to wait till the door was opened before she told \textit{her} story. She came hallooing to the window, <How do you do, my dear? How does Mrs Dashwood do? And where are your sisters? What! all alone! you will be glad of a little company to sit with you. I have brought my other son and daughter to see you. Only think of their coming so suddenly! I thought I heard a carriage last night, while we were drinking our tea, but it never entered my head that it could be them. I thought of nothing but whether it might not be Colonel Brandon come back again; so I said to Sir John, I do think I hear a carriage; perhaps it is Colonel Brandon come back again\longdash>

Elinor was obliged to turn from her, in the middle of her story, to receive the rest of the party; Lady Middleton introduced the two strangers; Mrs Dashwood and Margaret came down stairs at the same time, and they all sat down to look at one another, while Mrs Jennings continued her story as she walked through the passage into the parlour, attended by Sir John.

Mrs Palmer was several years younger than Lady Middleton, and totally unlike her in every respect. She was short and plump, had a very pretty face, and the finest expression of good humour in it that could possibly be. Her manners were by no means so elegant as her sister's, but they were much more prepossessing. She came in with a smile, smiled all the time of her visit, except when she laughed, and smiled when she went away. Her husband was a grave looking young man of five or six and twenty, with an air of more fashion and sense than his wife, but of less willingness to please or be pleased. He entered the room with a look of self-consequence, slightly bowed to the ladies, without speaking a word, and, after briefly surveying them and their apartments, took up a newspaper from the table, and continued to read it as long as he staid.

Mrs Palmer, on the contrary, who was strongly endowed by nature with a turn for being uniformly civil and happy, was hardly seated before her admiration of the parlour and every thing in it burst forth.

<Well! what a delightful room this is! I never saw anything so charming! Only think, Mama, how it is improved since I was here last! I always thought it such a sweet place, ma'am! (turning to Mrs Dashwood) but you have made it so charming! Only look, sister, how delightful every thing is! How I should like such a house for myself! Should not you, Mr Palmer?>

Mr Palmer made her no answer, and did not even raise his eyes from the newspaper.

<Mr Palmer does not hear me,> said she, laughing; <he never does sometimes. It is so ridiculous!>

This was quite a new idea to Mrs Dashwood; she had never been used to find wit in the inattention of any one, and could not help looking with surprise at them both.

Mrs Jennings, in the meantime, talked on as loud as she could, and continued her account of their surprise, the evening before, on seeing their friends, without ceasing till every thing was told. Mrs Palmer laughed heartily at the recollection of their astonishment, and every body agreed, two or three times over, that it had been quite an agreeable surprise.

<You may believe how glad we all were to see them,> added Mrs Jennings, leaning forward towards Elinor, and speaking in a low voice as if she meant to be heard by no one else, though they were seated on different sides of the room; <but, however, I can't help wishing they had not travelled quite so fast, nor made such a long journey of it, for they came all round by London upon account of some business, for you know (nodding significantly and pointing to her daughter) it was wrong in her situation. I wanted her to stay at home and rest this morning, but she would come with us; she longed so much to see you all!>

Mrs Palmer laughed, and said it would not do her any harm.

<She expects to be confined in February,> continued Mrs Jennings.

Lady Middleton could no longer endure such a conversation, and therefore exerted herself to ask Mr Palmer if there was any news in the paper.

<No, none at all,> he replied, and read on.

<Here comes Marianne,> cried Sir John. <Now, Palmer, you shall see a monstrous pretty girl.>

He immediately went into the passage, opened the front door, and ushered her in himself. Mrs Jennings asked her, as soon as she appeared, if she had not been to Allenham; and Mrs Palmer laughed so heartily at the question, as to show she understood it. Mr Palmer looked up on her entering the room, stared at her some minutes, and then returned to his newspaper. Mrs Palmer's eye was now caught by the drawings which hung round the room. She got up to examine them.

<Oh! dear, how beautiful these are! Well! how delightful! Do but look, mama, how sweet! I declare they are quite charming; I could look at them for ever.> And then sitting down again, she very soon forgot that there were any such things in the room.

% \begin{figure}[tbph]
% \centering
% \includegraphics[width=\linewidth]{19charming}
% \caption{<I declare they are quite charming>}
% \end{figure}

\begin{bwbigpic}
	[1.0]
	{19charming} 
	{<I declare they are quite charming>} 
\end{bwbigpic}

When Lady Middleton rose to go away, Mr Palmer rose also, laid down the newspaper, stretched himself and looked at them all around.

<My love, have you been asleep?> said his wife, laughing.

He made her no answer; and only observed, after again examining the room, that it was very low pitched, and that the ceiling was crooked. He then made his bow, and departed with the rest.

Sir John had been very urgent with them all to spend the next day at the park. Mrs Dashwood, who did not chuse to dine with them oftener than they dined at the cottage, absolutely refused on her own account; her daughters might do as they pleased. But they had no curiosity to see how Mr and Mrs Palmer ate their dinner, and no expectation of pleasure from them in any other way. They attempted, therefore, likewise, to excuse themselves; the weather was uncertain, and not likely to be good. But Sir John would not be satisfied—the carriage should be sent for them and they must come. Lady Middleton too, though she did not press their mother, pressed them. Mrs Jennings and Mrs Palmer joined their entreaties, all seemed equally anxious to avoid a family party; and the young ladies were obliged to yield.

<Why should they ask us?> said Marianne, as soon as they were gone. <The rent of this cottage is said to be low; but we have it on very hard terms, if we are to dine at the park whenever any one is staying either with them, or with us.>

<They mean no less to be civil and kind to us now,> said Elinor, <by these frequent invitations, than by those which we received from them a few weeks ago. The alteration is not in them, if their parties are grown tedious and dull. We must look for the change elsewhere.>